\section{6. Environment and
Constraints}\label{environment-and-constraints}

\subsection{6.1 Site and Physical
Environment}\label{site-and-physical-environment}

GSEG is an operational installation at one geographical site. Its design
must account for ambient temperature and humidity, precipitation and
wind, lightning and electromagnetic interference, fire and flood
hazards, and terrain constraints on antenna placement. Local values will
be established for the chosen site and equipment.

\subsection{6.2 Radio and Communications
Environment}\label{radio-and-communications-environment}

Ground-to-space communications depend on GEO/LEO visibility from the
site, atmospheric attenuation, RF interference, applicable spectrum
allocations and changes in available routes. SSEG provides direct radio
links and a mandatory ISL network; their existence does not itself
demonstrate link availability from the GSEG site.

\subsection{6.3 Infrastructure and Service
Availability}\label{infrastructure-and-service-availability}

GSEG depends on power, cooling, computing, networks, common time and
external communications provided or supported by GIS. A single site
creates common-cause failure exposure across mission functions. Two
on-site equipment domains do not eliminate complete-site hazards. GSEG
contributes to the EWAS \textbf{\textgreater99.9\%} system availability
requirement (\texttt{R-\allowbreak{}015}). Its separate critical-service operational
availability requirement is ≥99.99\%
(\texttt{GSEG-\allowbreak{}REQ-\allowbreak{}REL-\allowbreak{}003}), measured against the boundary in Section
3.2 and including common-site outages relevant to that boundary.

\subsection{6.4 Information Security}\label{information-security}

GSEG must protect spacecraft operational information, observations and
warnings through authorized access, protected communications, encryption
of archived data and traceable audit records, according to the relevant
system requirements. Detailed protection rules depend on information
classification and agreed interfaces.

\subsection{6.5 Operational Constraints}\label{operational-constraints}

The key constraints are the 60 s end-to-end warning deadline, the
two-LEO data prerequisite, variable spacecraft communication and data
availability, consistent timestamps and event identification, and
supervision of abnormal conditions. The first warning is automatic;
later operator confirmation, revision or cancellation is supported by
MCC/EWS.

\subsection{6.6 Regulatory and External
Conditions}\label{regulatory-and-external-conditions}

Operation is subject to applicable spectrum and antenna regulations,
facility safety, information-protection duties and arrangements with
external data providers and authorized warning recipients. Applicable
obligations depend on the selected site and external services. EWAS
warning responsibility ends at the authorized recipient interface
(\texttt{R-\allowbreak{}048}).
